\SourceLeaf{176}\par\Needspace{6\baselineskip}\medskip
\begin{center}\hypertarget{native-176-article129}{}\pdfbookmark[2]{129.}{article-129}129.\end{center}
Quinque methodi hactenus expositae protinus ad totidem alias viam sternunt,
quae ab illis eo tantum differunt, quod pro $x$ et $y$ loco distantiarum a
terra, inclinatio orbitae atque longitudo nodi ascendentis accipiuntur. Hae
igitur methodi novae ita se habent:

I. Determinantur ex $x$ et $y$ duobusque locis geocentricis extremis secundum
artt. 74, 75 longitudines heliocentricae in orbita radiique vectores, atque
hinc et ex temporibus respondentibus omnia reliqua elementa; ex his denique
locus geocentricus pro tempore observationis mediae, cuius differentiae a loco
observato in longitudine et latitudine ipsas $X$ et $Y$ suppeditabunt.

Quatuor reliquae methodi in eo conveniunt, quod e positione plani orbitae
locisque geocentricis omnes tres longitudines heliocentricae in orbita radiique
vectores respondentes computantur. Dein autem

II. elementa reliqua determinantur e duobus locis extremis tantum atque
temporibus respondentibus; secundum haec elementa calculantur pro tempore
observationis mediae longitudo in orbita atque radius vector, quarum quantitatum
differentiae a valoribus prius inventis, i. e. e loco geocentrico deductis,
ipsas $X,Y$ exhibebunt.

III. Aut derivantur orbitae dimensiones reliquae ex omnibus tribus locis
heliocentricis (artt. 82, 83), in quem calculum tempora non ingrediuntur: dein
temporum intervalla eruuntur, quae in orbita ita inventa inter observationem
primam et secundam, atque inter hanc et tertiam elapsa esse deberent, et
quorum differentiae a veris ipsas $X,Y$ nobis suggerent.

IV. Calculantur elementa reliqua duplici modo, puta tum e combinatione
loci primi cum secundo, tum e combinatione secundi cum tertio, adhibitis
temporum intervallis respondentibus. Comparatis hisce duobus elementorum
systematibus inter se, e differentiis duae quaecunque pro $X,Y$ accipi poterunt.

V. Sive denique idem calculus duplex tantummodo usque ad valores quantitatis
in art. 91 per $y$ denotatae producitur, ac dein pro $X,Y$ expressiones
in art. praec. traditae adoptantur.

Ut quatuor ultimis harum methodorum tuto uti liceat, loci terrae pro
omnibus tribus observationibus orbitae nodis non nimis vicini esse debent:
\SourceLeaf{177}contra usus methodi primae tamtummodo requirit, ut eadem conditio in duabus
observationibus extremis locum habeat, sive potius (quoniam locum medium
pro aliquo extremorum substituere licet), ut e tribus locis terrae non
plures quam unus in nodorum viciniis versentur.
